% =====================================================================
% 09 — Conclusion: where the hard part should live.
% =====================================================================
% Planning (% comments only):
% Tie everything back to the three-indictment partition.
% State what we will and will not claim until the companion suite reports.
% Articulate the future directions in concrete falsifiable terms.
% =====================================================================

\section{Where the Hard Part Should Live}
\label{sec:conclusion}

This paper organized the standard objections to autoregressive language models, separated them into three logically independent indictments, evaluated each against a clear bar, and surveyed the alternatives that target each.
A summary of the position the paper supports is below.

\paragraph{What survives, and what it indicts.}
Three families of objections survive the audit of \cref{sec:what_holds}.
The softmax bottleneck, doubly-exponential rank collapse in pure attention, and the local partition function tax indict the parameterization of the per-step distribution and the per-layer attention.
The mode-covering pressure of forward-KL training indicts the maximum-likelihood objective on noisy corpora.
The fixed-compute ceiling and the Lipschitz--margin gap indict the inference procedure of single-pass blind rollout.
The $(1-\varepsilon)^T$ argument is mathematically valid and rhetorically useful, but its three premises are explicitly violated by verifier-augmented systems and by feedback loops with hard checkers;
treating it as an impossibility proof confuses a warning label for a theorem.

\paragraph{Where the alternatives put the hard part.}
Energy-based models address the local-then-global mismatch by scoring whole objects.
The cost is the partition function or its approximation, paid either at training time (MLE with MCMC, contrastive divergence), at inference time (score-matching plus Langevin or diffusion sampling), or at neither but through a noise distribution whose coverage becomes the bottleneck (NCE).
Diffusion and flow-matching models share the inference-time iteration of score-matching;
they have largely succeeded in continuous data and are an active question in discrete data.
Joint-embedding predictive architectures avoid token-level reconstruction and the partition function but inherit the contrastive-versus-non-contrastive design space of self-supervised representation learning.
State-space models target only the per-layer partition function and, in their pure form, inherit the fixed-state recall limit;
hybrid stacks restore softmax attention at some scale and recover the property.
Autoregressive models with verifiers, search, retrieval, and tools attack the third indictment directly by negating the unrecoverability premise and the fixed-compute premise, and are the most empirically successful production response.

\paragraph{What this paper deliberately defers.}
This paper is the literature-grounded half of a two-part project.
The empirical claims --- which mechanisms are load-bearing inside a trained model on a domain with a hard verifier, and how the realized error process compares to the geometric prediction of \cref{ssec:compounding} --- are deferred to the companion suite, \texttt{experiments/02\_ar\_pros\_cons}, which pre-registers a probe-by-probe accounting at the correct layer for each surviving claim.
We will not commit to a verdict on the bridge between isolated mechanism strength and downstream realized effect until that suite reports.

\paragraph{What we are willing to say now.}
The honest summary of the literature is that the most-discussed critique of autoregressive language models, $(1-\varepsilon)^T$, is the weakest one in the surviving set, and the most-discussed alternative, ``replace softmax with an energy,'' is, by the Hopfield reading, a question of \emph{which} energy rather than \emph{whether} an energy.
The load-bearing research questions are narrower and more concrete than the slogans suggest:
\begin{enumerate}
  \item For each load-bearing isolated mechanism in \cref{tab:verdicts}, what is its realized effect inside an end-to-end trained model on a hard-verifier domain?
  \item For each candidate energy, which of the five load-bearing softmax properties does it preserve, and at what cost?
  \item Under what conditions is $K > 1$ steps of energy descent at inference cheaper than the equivalent additional training compute?
  \item What is the empirical decay law of $(1-\varepsilon_{\mathrm{effective}})^T$ for verifier-augmented generation on Lean, and is it geometric, sub-geometric, or polynomial?
\end{enumerate}
Each of these is a measurement question with a stated null and a stated alternative.
The companion suite is structured to answer them.
The answer will determine whether the architectural questions this paper raises remain open or are resolved by ablation.
